\section{Method}
\label{sec:methods}

3DPS renders complex-valued FMCW range profiles from a learned set of \emph{oriented 3D points} anchored on a LiDAR-derived geometric scaffold. Each point carries a quaternion-encoded surface normal and six physically interpretable ITU-R~P.2040 material parameters, evaluated by a closed-form bidirectional scattering distribution function (BSDF) and splatted into range-profile bins through a precomputed sensor point spread function (PSF). The pipeline is fully deterministic: there is no Monte Carlo sampling, no learned antenna pattern, and no random-direction integration. Sec.~\ref{sec:methods:fwd} states the radar forward model. Sec.~\ref{sec:methods:splat} presents the deterministic point-splatting renderer, our central contribution, and the LiDAR-derived geometric prior. Sec.~\ref{sec:methods:bsdf}--\ref{sec:methods:opt} specify the per-point ITU BSDF and the differentiable optimization.

% ---------------------------------------------------------------------
\subsection{Radar forward model}
\label{sec:methods:fwd}

For a single TX--RX pair observing a point target at $\mathbf{x}$ with normal $\mathbf{n}$, the bistatic radar equation~\cite{skolnik2001introduction} relates received power to bistatic radar cross-section, transmit power $P_t$, antenna gains $G_t,G_r$, wavelength $\lambda$, and TX/RX path lengths $d_t, d_r$. Replacing the point-target RCS with a spatially varying BSDF $f(\boldsymbol{\omega}_t, \mathbf{x}, \boldsymbol{\omega}_r)$ and integrating over the receiver hemisphere $\Omega_{rx}$~\cite{mmIR2026},
\begin{equation}
P_r \;=\; \frac{P_t\,\lambda^2}{(4\pi)^2}\!\!\int_{\Omega_{rx}}\!\!
\frac{G_t(\boldsymbol{\omega}_t)\,G_r(\boldsymbol{\omega}_r)\,
        f(\boldsymbol{\omega}_t,\mathbf{x},\boldsymbol{\omega}_r)\,c_t\,V}
       {d_t^2}\,d\omega_r,
\label{eq:radar_solid_angle}
\end{equation}
where $c_t = \max(0, \mathbf{n}\!\cdot\!\boldsymbol{\omega}_t)$ is the cosine of incidence and $V \in \{0,1\}$ is binary visibility. Standard FMCW renderers approximate Eq.~\eqref{eq:radar_solid_angle} via Monte Carlo path sampling, synthesize ADC samples per path, and apply a range FFT to recover range bins; mmIR~\cite{mmIR2026} is the current state of the art for this pipeline and requires ${\sim}$20~minutes for a 500-iteration single-viewpoint fit at 24k surface interactions per iteration.

% ---------------------------------------------------------------------
\subsection{Deterministic point splatting}
\label{sec:methods:splat}

\textbf{Why oriented 3D points, not Gaussians?} The radar equation in Eq.~\eqref{eq:radar_solid_angle} evaluates the BSDF at deterministic incidence and observation angles per (TX,~RX) pair, requiring explicit per-primitive position, surface normal, and material parameters at every scattering site. Oriented 3D points carry these parameters directly. The MC ray tracers we replace (mmIR~\cite{mmIR2026}, Sionna~RT~\cite{sionna}) already compute scene hit points per iteration to synthesize ADC; promoting those transient hits into persistent learnable primitives eliminates the per-iteration ray-tracing cost while preserving the physical structure that makes hit-based rendering correct in the first place. 3D Gaussians~\cite{3dgs} instead parameterize anisotropic density distributions with normals emergent from covariance eigenvectors and material properties learned implicitly via spherical harmonics, an indirect representation that does not align with the radar equation's requirement for explicit, deterministic per-point material and normal evaluation. The 3DGS-era arguments against point primitives in optics (``holes, aliasing, strict discontinuity''~\cite{3dgs}) calibrate to sub-pixel optical resolution and do not bind at the cascade radar's $\Delta r{=}5.93$~cm range and ${\sim}1.4^\circ$ azimuth resolution, where each oriented point splats a continuous Hann-windowed PSF spanning fifteen range bins (${\sim}89$~cm) and multiple azimuth bins after the virtual-array FFT.

\textbf{Deterministic evaluation.} 3DPS evaluates the forward model deterministically: each oriented point in the scene contributes \emph{exactly once} per (TX,~RX) pair, with no Monte Carlo path sampling and no random-direction integration. The hemisphere integral of Eq.~\eqref{eq:radar_solid_angle} is replaced by a closed-form sum over the fixed point set $\{\mathbf{x}_i\}_{i=1}^{N}$ with $N{=}20{,}000$, evaluated at the deterministic per-pair incidence and observation angles obtained directly from the array geometry.

\textbf{Range-profile splatting.} The range FFT of a single-path ADC contribution is a Hann-windowed sinc pulse centered at the path's bistatic range bin, fully determined by the chirp bandwidth and the FFT window. We precompute this point-spread function $\mathrm{PSF}(r) \in \mathbb{C}^K$ once on a $K{=}256$ bin grid at the radar's range resolution $\Delta r = c/(2\,B_{\text{chirp}}) = 5.93$~cm, where $B_{\text{chirp}}$ is the chirp's swept bandwidth. The forward pass then becomes a sparse scatter-add directly into range-profile bins:
\begin{equation}
\mathrm{RP}_{T R}[k]
\;=\;
\sum_{i=1}^{N}\,A_i^{T R}\,\exp(-j\,\phi_{T R i})\,
\mathrm{PSF}\!\bigl(k - k^{\star}_{T R i}\bigr),\quad
k^{\star}_{T R i} = \mathrm{round}(r_{T R i}/\Delta r),
\label{eq:psf_splat}
\end{equation}
where $A_i^{T R} \in \mathbb{C}$ collects the antenna gains, BSDF, and geometric falloff factors of Eq.~\eqref{eq:radar_solid_angle} for point $i$, $\phi_{T R i} = 2\pi(d_{T i}{+}d_{R i})/\lambda$ is the round-trip carrier phase, and $r_{T R i} = (d_{T i}{+}d_{R i})/2$ is the bistatic range. Splatting into the $L{=}15$ nearest range bins per path replaces the entire fast-time DFT, dropping the cost from $\mathcal{O}(N\!\cdot\!K)$ to $\mathcal{O}(N\!\cdot\!L)$ per (TX,~RX) pair.

\textbf{Validity of the splat.} The Hann-windowed PSF tail beyond $L{=}15$ taps is below $-50$~dB relative to peak, well under the noise floor of real ColoRadar captures. The bin-rounding $k^\star = \mathrm{round}(r/\Delta r)$ introduces at most $\Delta r/2 = 2.97$~cm of range bias per path, an order of magnitude below the radar's intrinsic range resolution; the rounding is applied only to the splat index and does not affect the carrier phase $\phi$, which is computed from the un-rounded $d_T+d_R$. Sidelobes and FFT windowing artifacts are captured exactly by the precomputed PSF. The splat is therefore numerically equivalent to the standard ADC-then-FFT pipeline on the same path set: we validate this empirically against mmIR's MC reference at maximum forward error $1.19\!\times\!10^{-7}$ on 90k paths~\cite{mmIR2026}.

\textbf{MIMO factorization.} The cascade radar's $N_{\text{TX}}{=}12$ transmitters and $N_{\text{RX}}{=}16$ receivers produce $192$ TX--RX pairs per chirp. The ITU BSDF (Sec.~\ref{sec:methods:bsdf}) factors over incidence and outgoing directions, so we evaluate the per-point incidence factor $f_i^{\text{in}}(\hat{\mathbf{d}}_{Ti})$ once per TX, the outgoing factor $f_i^{\text{out}}(\hat{\mathbf{d}}_{Ri})$ once per RX, and combine them per pair without recomputing material physics --- dropping the BSDF cost from $\mathcal{O}(N\!\cdot\!N_{\text{TX}}\!\cdot\!N_{\text{RX}})$ to $\mathcal{O}(N\,(N_{\text{TX}}{+}N_{\text{RX}}))$, the ray-reuse strategy of mmIR~\cite{mmIR2026,9533181}. After splatting, an $N_{\text{TX}}$-point virtual-array DFT along the azimuth dimension converts the per-(TX,~RX) range profiles into a range--azimuth (RA) map of shape $(N_{\text{az}},K) = (127, 256)$. The $N_{\text{az}}{=}127$ azimuth bins correspond to the unique TDM-MIMO virtual-element positions of the cascade array geometry, padded to the standard arcsin-spaced grid $\theta_a = \arcsin(2(a{-}N_{\text{az}}/2)/(N_{\text{az}}{+}1))$ with forward-cone half-angle $\arcsin(126/128) \approx 79.86^\circ$.

\textbf{CUDA fusion.} The combined cost of Eq.~\eqref{eq:psf_splat} plus factorized BSDF plus virtual-array DFT is implemented as three fused CUDA kernels: a forward BSDF kernel evaluating the GGX, SPM, and diffuse lobes plus Jones-polarized Fresnel coefficients in fp32; a scatter-splat kernel performing Eq.~\eqref{eq:psf_splat} via atomic adds into range-profile bins; and an analytical backward kernel computing per-point partial derivatives via atomic accumulation, with a closed-form Jacobian for the slab-Fresnel coefficients. The analytical backward is validated against finite-difference gradients to within numerical tolerance. The fused pipeline renders $N{=}20{,}000$ scene points across all $192$ (TX,~RX) pairs in ${\sim}10$~ms on an RTX~4090, and a 500-iteration training pass over the 8 training viewpoints completes in ${\sim}3.6$~minutes per scene.

\textbf{LiDAR-derived geometric prior.} We initialize the point set $\{\mathbf{x}_i\}$ from a co-registered LiDAR scan. A 3D geometric prior is mandatory, not optional: the forward model in Eqs.~\eqref{eq:radar_solid_angle}--\eqref{eq:psf_splat} requires full 3D positions $\mathbf{x}_i\in\mathbb{R}^3$ to compute bistatic range $r_{TRi}$, incidence/observation angles $\boldsymbol{\omega}_t,\boldsymbol{\omega}_r$, and the elevation-dependent antenna gains $G_t G_r$ (the cascade's planar virtual array has a non-trivial elevation pattern even though it lacks an elevation aperture). The supervision signal --- 2D range--azimuth maps --- has integrated out the elevation dimension by construction, so RA maps alone are information-theoretically insufficient to recover 3D point positions regardless of how many viewpoints are aggregated. Compounding this, NVS is performed from only eight training views per scene, and each individual RA map is itself extremely sparse: mmWave returns are dominated by strong specular reflections at near-normal incidence with sparse glints elsewhere, since outdoor surfaces tend to act as quasi-mirrors at 60--77\,GHz~\cite{guan2020through, lu2013measurement}. LiDAR-conditioned scaffolds are accordingly the dominant precedent for radar rendering and ray-tracing on real scenes~\cite{dudek2010millimeter, 9533181, mostajabi2020high, kramer2022coloradar, mmIR2026, kung2025radarsplat, rafidashti2025neuradar, sionna}. We reduce the per-scene LiDAR cloud (positions, normals, intensity; $\sim 2.7$--$6.5$M points across our scenes) to $N{=}20{,}000$ oriented points by (1) culling outside the radar's azimuth cone ($\cos(\angle(\hat{\mathbf{d}},\hat{\mathbf{b}})) > 0.1761$) and radial range ($1.5$~m to $K\Delta r{=}15.2$~m); (2) cosine-weighted resampling of $3N$ candidates with replacement at probability $\propto \max(\cos(\angle(\hat{\mathbf{d}},\hat{\mathbf{b}})), 0.01)$, biasing selection toward boresight-aligned regions where the antenna pattern is strongest; and (3) farthest-point sampling to exactly $N{=}20{,}000$. Each surviving point inherits a quaternion encoding its LiDAR surface normal and a uniform ITU-concrete material prior.

% Complexity comparison table (commented out for now; complexity discussion stays in prose above)
% \begin{table}[t]
\centering
\caption{Per-component comparison of 3DPS deterministic point splatting vs.\ mmIR's MC path tracer. $N$\,=\,scene point count, $K\,{=}\,256$ range bins, $L\,{=}\,15$ PSF taps per point, $N_{\text{TX}}\!\cdot\!N_{\text{RX}}\,{=}\,192$ TX--RX pairs. Wall-clock figures: 500 iterations on a single RTX~4090.}
\label{tab:complexity}
\small
\setlength{\tabcolsep}{6pt}
\renewcommand{\arraystretch}{1.15}
\begin{tabular}{lll}
\toprule
\textbf{Component} & \textbf{mmIR~\cite{mmIR2026}} & \textbf{3DPS (Ours)} \\
\midrule
Path sampling      & MC over $24$k paths / iter            & deterministic, $N{=}20$k points / iter (no MC) \\
Range processing   & ADC synth.\ + range FFT, $\mathcal{O}(N\!\cdot\!K)$ & PSF splat, $\mathcal{O}(N\!\cdot\!L)$ \\
BSDF evaluation    & per-pair, factorized over $(N_{\text{TX}}{+}N_{\text{RX}})$ & per-pair, factorized over $(N_{\text{TX}}{+}N_{\text{RX}})$ \\
Backward pass      & autograd through MC graph           & fused analytical CUDA kernel \\
\midrule
Forward / iter, 8 views, $N{=}20$k & ---  & ${\sim}10$~ms \\
Wall-clock, 500 it., 1 view        & ${\sim}20$~min  & ---  \\
Wall-clock, 500 it., 8 views       & not reported (multi-view intractable) & $\boldsymbol{\sim}\textbf{2.3 min}$ (${\sim}70\times$ faster per viewpoint) \\
\bottomrule
\end{tabular}
\end{table}


\begin{figure}[t]
  \centering
  \includegraphics[width=\linewidth]{figs_highres/pipeline_3dps.pdf}
  \caption{\textbf{3DPS rendering pipeline.} A LiDAR point cloud is reduced to $N{=}20{,}000$ oriented points (Sec.~\ref{sec:methods:splat}, LiDAR-derived geometric prior). For each (TX,~RX) pair, the per-point ITU BSDF (Sec.~\ref{sec:methods:bsdf}) is evaluated at the deterministic incidence and observation angles, producing a complex amplitude $A_i^{TR}$ and a carrier phase $\phi_{TRi}$. Each point's complex contribution is splatted via a Hann-windowed PSF into the $L{=}15$ nearest range bins (Eq.~\ref{eq:psf_splat}). Per-(TX,~RX) range profiles are stacked and an azimuth FFT is applied to obtain the final RA map.}
  \label{fig:pipeline}
\end{figure}

% ---------------------------------------------------------------------
\subsection{Per-point ITU BSDF}
\label{sec:methods:bsdf}

Each point's BSDF follows the ITU-R~P.2040 multilayer slab model used by mmIR~\cite{mmIR2026,itu_p2040,wei2024csvbsdf}, applied per-point rather than per-vertex. The 6-parameter material vector $\boldsymbol{\theta}_i = (\varepsilon_r',\,\varepsilon_r'',\,\sigma_h,\,l_c,\,\tau,\,d)_i$ specifies the complex relative permittivity $\varepsilon_r = \varepsilon_r' - j\,\varepsilon_r''$, RMS surface roughness $\sigma_h$, correlation length $l_c$, Kirchhoff-approximation / small-perturbation (KA/SPM) blend $\tau \in [0,1]$, and slab thickness $d$. Each parameter is reparameterized via sigmoid (KA/SPM blend) or shifted softplus (positive-only quantities) to keep optimization in physically valid ranges.

The BSDF decomposes into coherent and incoherent components via the Rayleigh roughness factor $\eta = \exp(-(2k\sigma_h\cos\theta_t)^2)$ with $k = 2\pi/\lambda$:
\begin{equation}
f \;=\; A(\boldsymbol{\omega}_t)\,\bigl[\eta\,f_{\text{coh}} + (1-\eta)\,f_{\text{inc}}\bigr],
\label{eq:bsdf_decomp}
\end{equation}
where $A(\boldsymbol{\omega}_t) \in [0,1]$ is the slab-Fresnel energy gate (multilayer slab with thickness $d$ and complex permittivity $\varepsilon_r$, capturing the internal reflections that single-interface Fresnel misses --- critical for plasterboard, vehicle bodies, and other thin layered materials), $f_{\text{coh}} = \tau f_{\text{KA}} + (1-\tau) f_{\text{SPM}}$ combines a GGX microfacet specular lobe and a von~Mises--Fisher SPM lobe, and $f_{\text{inc}} = \gamma f_{\text{dir}} + (1-\gamma) f_{\text{broad}}$ combines a directional and a Lambertian diffuse lobe. Polarization is tracked through the local $s/p$ basis with complex Jones reflection coefficients per polarization. We reuse mmIR's per-lobe forms unchanged~\cite{mmIR2026}; per-point optimization yields one $\boldsymbol{\theta}_i$ vector per scene point rather than per mesh vertex.

% ---------------------------------------------------------------------
\subsection{Differentiable optimization}
\label{sec:methods:opt}

\textbf{Learned parameters.} For each point we learn (i) the 6-parameter ITU material $\boldsymbol{\theta}_i$ (shared init, per-point freedom afterward); (ii) the surface-normal quaternion $\mathbf{q}_i$, initialized from the LiDAR normal; and (iii) the 3D position $\mathbf{p}_i$, optimized through the \emph{amplitude path only}. The carrier-phase term $e^{-j\phi_{T R i}}$ in Eq.~\eqref{eq:psf_splat} is detached from the autograd graph for position gradients because the rounded range-bin index $k^\star = \mathrm{round}(r/\Delta r)$ creates a piecewise-constant dependence on $\mathbf{p}_i$: $\partial k^\star / \partial \mathbf{p}_i$ is zero almost everywhere and an unbounded delta at bin boundaries, so phase gradients on position are degenerate. Detaching the carrier phase removes the discontinuity from the backward pass without affecting the forward pass; the carrier phase is still computed from the un-rounded $d_T + d_R$. An L2 anchor $\lambda_{\text{pos}}\|\mathbf{p}_i - \mathbf{p}_i^{(0)}\|^2$ ($\lambda_{\text{pos}}{=}100$) restricts position drift to ${\sim}1$~mm per iteration --- much smaller than the carrier wavelength ($\lambda \approx 3.9$~mm at 76\,GHz) --- so geometry remains LiDAR-anchored and the carrier phase remains correct in expectation.

\textbf{Loss.} For each training pose $F$ we render the magnitude RA map $|\mathrm{RA}_F^{\text{pred}}|$ via Eq.~\eqref{eq:psf_splat} plus azimuth DFT and compute MSE against the measured $|\mathrm{RA}_F^{\text{GT}}|$, both normalized by their respective $L^2$ norms. The total loss is the mean over training samples; rotation and material drift are unregularized and bounded implicitly by Adam.

\textbf{Adaptive density control.} Following 3DGS~\cite{3dgs} but adapted to radar physics, we periodically split high-importance points and prune low-importance points to redistribute the fixed $N$ budget. The per-point importance signal is the accumulated amplitude-path position-gradient magnitude $g_i = \sum_F \|\nabla_{\mathbf{p}_i} \mathcal{L}_F\|$ between densify rounds: a high $g_i$ indicates the renderer disagrees with the data at that point, so the point either needs additional capacity (split) or is contributing nothing (prune). To compute $g_i$ without actually moving points, we track the gradient on positions while \emph{excluding} them from the optimizer step, so positions remain on the LiDAR scaffold while the gradient drives the splitting decision. At iters $\{100, 200, 300, 400\}$ we (i) clone the top 5\% of points by $g_i$ into Gaussian-jittered children at $\sigma_p = 2$~cm, $\sigma_q = 0.05$~rad, (ii) evict the bottom 5\% by $g_i$, and (iii) reset Adam state for the affected slots. The budget $N{=}20{,}000$ is exactly preserved.

\textbf{Implementation.} Optimization uses Adam ($\beta_1{=}0.9$, $\beta_2{=}0.999$) with per-parameter learning rates $\eta_{\text{mat}}{=}10^{-2}$, $\eta_{\text{rot}}{=}5\!\cdot\!10^{-3}$, $\eta_{\text{pos}}{=}10^{-5}$, and per-parameter root-mean-square gradient clipping. Each scene's 500-iteration training loop completes in ${\sim}3.6$~minutes on a single RTX~4090.
